\chapter{Foundations and Transformative Scope of \hydrot}

\phantomsection
\fancyhead[LE]{\thepage} 
\fancyhead[RE]{\textbf{Hydrological Twin - Foundations and Transformative Scope}}
\fancyhead[RO]{\thepage} 
\fancyhead[LO]{\textbf{Hydrological Twin - Foundations and Transformative Scope}}
\renewcommand{\headrulewidth}{1pt}

\setlength{\parindent}{0cm} % no indent
\setlength{\parskip}{9pt}

% --------------------------------------------------------------


\section{\hydrot\ in brief}


\hydrot\ establishes a new epistemic regime for hydrological modelling. It replaces loosely coupled numerical experiments by synchronized, stateful, traceable and reproducible anchored estimations. It is multi--dataset and multi--model by construction, natively compatible with Bayesian and AI methodologies, and structurally committed to audit given a native open--source design.

By isolating physical model parameterizations as the only potentially private component, \hydrot\ reconciles scientific openness with operational and economic constraints. In doing so, it provides a rigorous and future-proof foundation for digital hydrology, at the intersection of physics, data science, and distributed versioned infrastructures.

\hydrot\ is a digital twin of hydrological systems implemented as a Python package following a monolithic architectural approach. It is conceived as a persistent artefact evolving in the ontological space-state continuum.








Each \hydrot\ instance is linked to a complex ecosystem of Git repositories encompassing databases, models, and configuration artefacts, each of them being carefully registered and versioned in terms of space-state.





At its core, \hydrot\ relies on a small number of conceptual classes that formalize identity, versioning, spatial supports, observations, simulations, and model configurations. In particular, identity and traceability are enforced through a \texttt{GitSynchroRecord} that captures repository provenance (commit, branch, tag/describe, dirty state, and remote URL) and uses it to compute a reproducible version fingerprint through deterministic Git provenance and environment fingerprinting.


\section{Architecture}

One instance of a \hydrot\ is composed of six functional layers :

\begin{enumerate}
    \item Model Layer (process-based physical models -- software, spatial and temporal supports of the variable - grid)
    \item Data Layer (observations, reanalyses, forcings, inner parameters - private)
    \item Estimation Layer (comparison, filtering, Bayesian inference)
    \item Analysis Layer (temporal and spatial transformations, extraction - on demand)
    \item Cartographic Layer (visualization and spatial representation)
    \item Git-Synchronized Registry (identity, provenance, versioning)
\end{enumerate}

\section{Foundational Principles}

\hydrot\ is conceived as a foundational digital infrastructure for hydrological science and engineering. Its primary objective is not merely to host numerical models, but to formalize hydrological knowledge as a persistent, traceable, and evolutive artefact in the space-state continuum.

The core transformation introduced by \hydrot\ is the passage from classical physically--based modelling toward \emph{synchronized physically--based estimation}. In conventional practice, physically--based models are executed as isolated numerical experiments, producing results that are weakly linked to their exact computational, parametric, and data provenance. In contrast, \hydrot\ enforces that every simulation, dataset, parameterization, and software component is explicitly registered, versioned, and fingerprinted within a coherent space-state framework.

This shift elevates hydrological modelling from a sequence of ad hoc numerical runs to a formally structured process of stateful estimation, where each computational act becomes a reproducible, auditable, and referencable object.

\subsection{From Physically--Based Models to Physically--Based Estimation}

The defining transformative property of \hydrot\ lies in its ability to synchronize physically--based models with the evolving space-state of the digital twin.

In this paradigm, a hydrological model is no longer treated as a static algorithmic object, but as a dynamic estimation engine whose outputs are strictly bound to:
\begin{itemize}
\item a precise software state (source code, numerical kernels, compilation environment),
\item an explicit set of forcings and observational products,
\item a uniquely identified parameterization,
\item a well-defined spatial support,
\item and a fully captured Git provenance snapshot.
\end{itemize}

Each execution therefore constitutes a physically--based estimation step embedded in a synchronized space-state. Any modification of code, data, parameters, or repository state deterministically generates a new version fingerprint, making silent drift or hidden inconsistencies structurally impossible.

This design transforms model outputs into first-class scientific objects: they are no longer merely numbers, but estimators whose identity and lineage are cryptographically anchored to their full computational context.

\subsection{Multi--Dataset and Multi--Model by Construction}

\hydrot\ is multi--dataset and multi--model by construction. This property does not arise from a plug-in layer added a posteriori, but from the very structure of its core object model.

Multiple observational products, forcings, and databases can coexist within the same digital twin instance, each one being registered through its own identity card and Git provenance. Likewise, several physically--based models or numerical engines can be bound to the same hydrological system through distinct \texttt{Software} and \texttt{Simulation} objects.

As a consequence, \hydrot\ naturally supports:
\begin{itemize}
\item ensemble modelling across structurally different models,
\item intercomparison of alternative datasets or reanalyses,
\item hybrid modelling strategies combining complementary physical formulations,
\item progressive refinement of spatial supports and discretizations.
\end{itemize}

These capabilities emerge without introducing conceptual or technical tension, because the digital twin is not tied to a single canonical model or dataset. Instead, it formalizes a structured coexistence of multiple, versioned representations of the same hydrological reality.

\subsection{Ontological State-Space of the Hydrological Twin}

The Hydrological Twin framework is explicitly grounded in the notion of an \emph{ontological state-space} for hydrological simulations. This concept departs from the traditional view of simulations as isolated model runs and instead formalizes each Hydrological Twin instance as a realization of a structured state-space whose dimensions, semantics, and admissible transformations are explicitly defined.

Within this framework, a Hydrological Twin is not merely the numerical output of a model executed under a given set of forcings and parameters. It is a well-defined point, or trajectory, in a high-dimensional state-space that integrates physical variables, forcings, parameters, and diagnostics across multiple spatial and temporal scales. The structure of this space is explicit, documented, and shared, enabling reproducibility, comparison, and composability across modeling chains.

Each state variable embedded in the Hydrological Twin is treated as an \emph{ontological object}. Beyond its numerical value, it is characterized by a set of attributes that define its meaning and validity. These attributes include: \begin{itemize}
    \item its provenance (e.g.\ originating global climate model, emissions scenario, downscaling or filtering methodology),
    \item its spatial scale (global, regional, or local)
    \item its functional role within the system (forcing, prognostic state, parameter, or diagnostic variable)
    \item its domain of validity in time, space
    \item underlying physical or modeling assumptions. 
 \end{itemize} 
 This explicit typing prevents semantic ambiguity when combining heterogeneous datasets and models.

Transitions between applications within the Hydrological Twin ecosystem—such as the progression from global-scale climate data aggregation to regional downscaling and finally to hydrosystem simulation—are formalized as \emph{morphisms between ontological sub-spaces}. In this view, scale change is not an ad-hoc numerical operation but a well-defined transformation that maps one structured state-space onto another while preserving documented semantics and constraints. Interpolation, bias correction, and downscaling are therefore interpreted as specific classes of state-space transformations rather than isolated preprocessing steps.

Estimation procedures and data assimilation are naturally embedded within this ontological framework. They are reinterpreted as processes of \emph{state synchronization} within the state-space, whereby observational information constrains and updates physically based model states. This perspective aligns with the concept of physically based estimation synchronized in state-space and avoids treating assimilation as an external numerical correction detached from the model ontology.

Finally, version control plays a foundational role in the governance of the ontological state-space. Git-based versioning is not limited to tracking source code evolution; each committed Hydrological Twin corresponds to a uniquely identifiable point or trajectory in the ontological state-space. This enables full traceability of simulations, including their structural definitions, transformations, and parameterizations, and allows systematic comparison between alternative twins. This property fundamentally distinguishes the Hydrological Twin from standard digital twin implementations, which typically lack formal state-space semantics and long-term traceability.

% --- AJOUT DANS 1.3.3 (Ontological State-Space) ---
Within the ontological state-space, \textbf{private-key authenticated users} can:
\begin{itemize}
    \item Access raw state variables and full-resolution data.
    \item Modify parameterizations and execute new simulations.
\end{itemize}
Public users interact with a \textbf{restricted state-space subset}, limited to aggregated and/or masked variables.

% --- FIN DE L'AJOUT ---



\subsection{Multi simulation by design : Space--State as a Substrate for Bayesian and AI Methods}

The explicit multi simulation space-state formalism composed of multiple simulations implemented in \hydrot\ provides a natural substrate for Bayesian inference and for the deployment of artificial intelligence methods.

Because each simulation and dataset is uniquely identified and versioned, posterior distributions, likelihood evaluations, surrogate models, and learned emulators can be attached to precise space-state coordinates. This allows Bayesian updating schemes in which priors, posteriors, and evidence are not abstract mathematical objects, but concrete, traceable artefacts stored within the digital twin.

In this context, AI models are no longer black boxes trained on loosely documented datasets. They become estimators whose training data, architecture, hyperparameters, and resulting weights are themselves registered as space-state objects with full Git provenance.

This design enables:
\begin{itemize}
\item rigorous reproducibility of machine learning workflows,
\item controlled hybridization between physical and data-driven models,
\item auditable deployment of AI in operational hydrology,
\item Bayesian model averaging and data assimilation directly grounded in the digital twin.
\end{itemize}

\subsection{Open--Source, Traceable, and Service--Oriented by Design}

\hydrot\ is conceived by construction as an open--source, service--oriented infrastructure. Its codebase, datasets, and configuration artefacts are intended to be hosted in publicly accessible Git repositories whenever legally and ethically possible.

Traceability is not an optional feature but a structural invariant: every \hydrot\ object carries a \texttt{GitSynchroRecord} embedding commit hashes, branch or tag identifiers, dirty state, and remote URLs. The universal registry further anchors each validated digital twin instance in an append-only, hash-chained ledger synchronized on Git.

This architecture turns Git into a scientific instrumentation layer, providing cryptographic integrity, distributed synchronization, and long-term archival capabilities at negligible marginal cost.

Every object that participates in the twin space-state must expose a \texttt{GitSynchroRecord}, removing symbolic identity in favor of computational ontology.

% --- AJOUT DANS 1.3.5 (Open-Source, Traceable, and Service-Oriented) ---
Traceability is enforced for all actions: public operations are openly logged, while \textbf{private-key authenticated actions} (e.g., data extraction, parameter changes) are tied to user identities and cryptographically recorded. This dual-layer approach ensures both transparency and confidentiality where required.

% --- FIN DE L'AJOUT ---



% --- NOUVELLE VERSION DE 1.3.6 ---
\subsection{The Only Private Components: Instanciation, Physical Model Parameterizations and Extraction from Simulation Hyper-Cube}
\label{sec:private-components}

By design, HydrologicalTwin isolates \textbf{two categories of private components}:
\begin{enumerate}
    \item \textbf{Instanciation} TODO 
    \item \textbf{Physical model parameterizations}: These may encode strategic or proprietary knowledge and are stored in restricted-access repositories. Their existence and version are traceable via \texttt{GitSynchroRecord}, but their numerical content remains confidential.
    \item \textbf{Data extraction and analysis operations}: Raw simulation data (e.g., full-resolution NumPy hypercubes) and advanced analytical operations (e.g., \texttt{extract\_values}, \texttt{apply\_temporal\_operator}) require \textbf{private-key authentication}. Without a valid key, users are limited to:
    \begin{itemize}
        \item Masked or aggregated outputs (e.g., basin-level averages).
        \item Visualization layers (e.g., spatial maps, time-series plots) via the \textit{Rendering} layer.
    \end{itemize}
\end{enumerate}

\subsubsection*{Access Control Mechanism}
\begin{itemize}
    \item Public users interact with the twin through a \textbf{restricted API}, accessing only pre-computed or masked data.
    \item Private-key holders unlock full functionality, including:
    \begin{itemize}
        \item Raw data extraction.
        \item Parameter modification.
        \item Execution of new simulations.
    \end{itemize}
    \item All private actions are \textbf{cryptographically logged} to ensure traceability and prevent tampering.
\end{itemize}

\subsubsection*{Rationale}
This separation reconciles scientific openness with operational constraints:
\begin{itemize}
    \item \textbf{Transparency}: Public users can verify the existence and provenance of simulations via \texttt{GitSynchroRecord}.
    \item \textbf{Confidentiality}: Sensitive parameters and raw data remain protected.
    \item \textbf{Reproducibility}: Private actions are tied to user identities via key-based authentication, ensuring accountability.
\end{itemize}

% --- FIN DE LA SECTION ---
